\documentclass[12pt,a4paper]{article}

\usepackage[utf8]{inputenc}
\usepackage{geometry}
\geometry{a4paper, margin=1in}

\usepackage{graphicx}
\usepackage{float}
\usepackage{hyperref}
\usepackage{booktabs}
\usepackage{xcolor}
\usepackage{caption}
\usepackage{subcaption}
\usepackage{amsmath}

\title{
    \vspace*{1.2cm}
    \Huge \textbf{Project State Report} \\
    \vspace{0.4cm}
    \LARGE \textit{Unified Multi-Domain Industrial Inspection Platform} \\
    \vspace{0.8cm}
}

\author{
    \Large \textbf{Done by:} Youssef Mahdi \\
    \vspace{0.3cm}
    \Large \textbf{Supervised by:} Zakaria Abouhanifa \& Rime Bouali \\
    \vspace{0.6cm}
    \Large \textbf{Company:} Nextronic
}

\date{\today}

\begin{document}

\begin{titlepage}
    \maketitle
    \thispagestyle{empty}
    \vfill
    \begin{center}
        \large \textit{Architecture, post-processing logic, and model-performance analysis for Stator, Chignon, and File domains.}
    \end{center}
\end{titlepage}

\tableofcontents
\newpage

\section{Executive Summary}

The application now runs as a \textbf{unified multi-domain inference system} with independent runtime paths for:
\begin{itemize}
    \item \textbf{Stator} detection,
    \item \textbf{Chignon} detection,
    \item \textbf{File} detection.
\end{itemize}

The infrastructure for all domains is integrated in one Flask + WebRTC application. Model behavior and post-processing are domain-specific.

\textbf{Critical model decision (File domain):}
YOLO file-domain runs are non-convergent in current training outputs (mask mAP metrics collapse to 0.0), so deployment uses \textbf{UNet-ResNet18 (.pth)} for file detection.

\section{Application Architecture}

\subsection{Domain-Aware Runtime}

Requests are routed by domain to dedicated managers:
\begin{itemize}
    \item Stator $\rightarrow$ \texttt{ModelManager}
    \item Chignon $\rightarrow$ \texttt{ChignonModelManager}
    \item File $\rightarrow$ \texttt{FileModelManager}
\end{itemize}

This routing is applied consistently to model listing, model switching, enhancement settings, upload detection, and WebRTC live inference.

\begin{figure}[H]
    \centering
    \includegraphics[width=0.58\textwidth,height=0.30\textheight,keepaspectratio]{docs/images/multi_domain_architecture.png}
    \caption{Multi-domain architecture: one app, three domain managers, domain-specific post-processing.}
\end{figure}

This architecture shows the core orchestration principle of the application: a single UI and transport layer delegates inference to independent domain managers. Each manager owns model discovery, model loading, prediction, and post-processing semantics for its domain. This separation avoids cross-domain rule leakage (for example, stator geometry rules affecting file logic) while preserving one deployment surface for APIs, WebRTC streaming, and settings.

\begin{figure}[H]
    \centering
    \includegraphics[width=0.58\textwidth,height=0.30\textheight,keepaspectratio]{docs/images/detection_pipeline.png}
    \caption{Unified detection flow with domain manager selection stage.}
\end{figure}

The unified flow emphasizes that preprocessing remains shared, but decision logic diverges right after domain selection. Input frames are normalized once, then routed to the selected manager where detection outputs are transformed by domain-specific post-processing. This design keeps latency predictable and allows independent evolution of stator measurement, chignon segmentation tuning, and file color-layout validation without changing the transport pipeline.

\subsection{Separate Detection Pipelines}

Each domain has its own processing flow due to different industrial constraints.

\begin{figure}[H]
    \centering
    \includegraphics[width=0.56\textwidth,height=0.23\textheight,keepaspectratio]{docs/images/stator_detection_pipeline.png}
    \caption{Stator pipeline: top-N filtering, spatial correction, contour geometry, and measurement calibration.}
\end{figure}

The stator pipeline is measurement-centric. After raw detections, top-N constraints enforce physically valid object counts, then a spatial heuristic corrects class confusion using position relative to the stator center. Contours are converted to geometric primitives, and calibrated scale factors convert pixel distances into millimeter-level outputs for QA metrics such as width, diameter, area, and inter-part spacing.

\begin{figure}[H]
    \centering
    \includegraphics[width=0.56\textwidth,height=0.23\textheight,keepaspectratio]{docs/images/chignon_detection_pipeline.png}
    \caption{Chignon pipeline: domain-specific segmentation and rendering with optional tracking/edge refinement hooks.}
\end{figure}

The chignon pipeline prioritizes stable segmentation and visualization fidelity over heavy geometric measurement. Domain-specific rendering logic controls class overlays and confidence presentation, while optional tracking/edge-refinement hooks can be enabled to smooth temporal behavior and contour boundaries. This keeps the path lightweight and suited for real-time monitoring with minimal rule-based overrides.

\begin{figure}[H]
    \centering
    \includegraphics[width=0.56\textwidth,height=0.23\textheight,keepaspectratio]{docs/images/file_detection_pipeline.png}
    \caption{File pipeline: ResNet inference, color classification, left-right validation, and correctness feedback.}
\end{figure}

The file pipeline is validation-centric and currently anchored on UNet-ResNet18 outputs. Detected file regions are post-processed by color scoring inside bounding boxes to classify blue/yellow semantics, then a left-right rule checks assembly positioning. The system draws a center line between paired files and emits explicit correctness messages to the UI, enabling operators to interpret not only detection presence but also placement conformity.

\section{Model Results Analysis by Domain}

\subsection{Stator Domain Results}

Training artifacts analyzed from \texttt{outputs/results}.

\begin{table}[H]
    \centering
    \begin{tabular}{lcccc}
        \toprule
        \textbf{Model} & \textbf{Best Epoch} & \textbf{Primary Metric} & \textbf{Value} & \textbf{Peak GPU (MB)} \\
        \midrule
        UNet-ResNet18 & 30 & Val IoU & 0.9663 & 472.5 \\
        YOLOv8m-seg & 30 & Mask mAP50-95 & 0.8751 & -- \\
        YOLOv11m-seg & 30 & Mask mAP50-95 & 0.8682 & -- \\
        \bottomrule
    \end{tabular}
    \caption{Stator model summary (from training logs and YOLO results CSV).}
\end{table}

Interpretation:
\begin{itemize}
    \item UNet-ResNet18 provides strongest pixel-level IoU for stator segmentation.
    \item YOLO models maintain high precision/recall and strong instance-level mask mAP.
    \item Stator domain supports both YOLO and PyTorch runtime paths in production.
\end{itemize}

\begin{figure}[H]
    \centering
    \begin{subfigure}{0.48\textwidth}
        \centering
        \includegraphics[width=\textwidth]{outputs/results/plots/iou_curves.png}
        \caption{Stator IoU curves}
    \end{subfigure}
    \hfill
    \begin{subfigure}{0.48\textwidth}
        \centering
        \includegraphics[width=\textwidth]{outputs/results/plots/loss_curves.png}
        \caption{Stator loss curves}
    \end{subfigure}
    \caption{Stator training dynamics.}
\end{figure}

\subsection{Chignon Domain Results}

Training artifacts analyzed from \texttt{chignon/results}.

\begin{table}[H]
    \centering
    \begin{tabular}{lcccc}
        \toprule
        \textbf{Model} & \textbf{Best Epoch} & \textbf{Primary Metric} & \textbf{Value} & \textbf{Peak GPU (MB)} \\
        \midrule
        UNet-ResNet18 & 39 & Val IoU & 0.9816 & 440.9 \\
        YOLOv8m-seg & 40 & Mask mAP50-95 & 0.9688 & -- \\
        YOLOv11m-seg & 39 & Mask mAP50-95 & 0.9606 & -- \\
        \bottomrule
    \end{tabular}
    \caption{Chignon model summary (from training logs and YOLO results CSV).}
\end{table}

Interpretation:
\begin{itemize}
    \item Chignon domain converges strongly for both UNet and YOLO variants.
    \item YOLO recall is saturated at 1.0 in best checkpoints, indicating high object recovery on this dataset.
    \item Domain is integrated in app routing and remains suitable for YOLO/PyTorch comparison workflows.
\end{itemize}

\begin{figure}[H]
    \centering
    \begin{subfigure}{0.48\textwidth}
        \centering
        \includegraphics[width=\textwidth]{chignon/results/plots/iou_curves.png}
        \caption{Chignon IoU curves}
    \end{subfigure}
    \hfill
    \begin{subfigure}{0.48\textwidth}
        \centering
        \includegraphics[width=\textwidth]{chignon/results/plots/loss_curves.png}
        \caption{Chignon loss curves}
    \end{subfigure}
    \caption{Chignon training dynamics.}
\end{figure}

\subsection{File Domain Results}

Training artifacts analyzed from \texttt{files/results}.

\begin{table}[H]
    \centering
    \begin{tabular}{lcccc}
        \toprule
        \textbf{Model} & \textbf{Best Epoch} & \textbf{Primary Metric} & \textbf{Value} & \textbf{Peak GPU (MB)} \\
        \midrule
        UNet-ResNet18 & 30 & Val IoU & 0.9197 & 440.9 \\
        YOLOv8m-seg & 1 & Mask mAP50-95 & 0.0000 & -- \\
        YOLOv11m-seg & 1 & Mask mAP50-95 & 0.0000 & -- \\
        \bottomrule
    \end{tabular}
    \caption{File model summary (from training logs and YOLO results CSV).}
\end{table}

Interpretation:
\begin{itemize}
    \item File-domain YOLO training failed to produce deployable mask quality (all key mask metrics collapse to zero).
    \item UNet-ResNet18 remains the only viable model path for file-domain inference in this cycle.
    \item Deployment therefore prioritizes UNet-ResNet18 plus file-specific post-processing rules.
\end{itemize}

\begin{figure}[H]
    \centering
    \begin{subfigure}{0.48\textwidth}
        \centering
        \includegraphics[width=\textwidth]{files/results/plots/iou_curves.png}
        \caption{File IoU curves}
    \end{subfigure}
    \hfill
    \begin{subfigure}{0.48\textwidth}
        \centering
        \includegraphics[width=\textwidth]{files/results/plots/loss_curves.png}
        \caption{File loss curves}
    \end{subfigure}
    \caption{File-domain training dynamics.}
\end{figure}

\section{Post-Processing Methods by Domain}

\subsection{Stator Post-Processing}
\begin{itemize}
    \item Top-N filtering limits structurally valid object counts.
    \item Spatial heuristic correction mitigates class confusion based on angular position around stator center.
    \item Contour extraction + geometry fitting computes dimensions.
    \item Calibration methods (manual/reference/intrinsics/depth) map px to mm.
\end{itemize}

\subsection{Chignon Post-Processing}
\begin{itemize}
    \item Domain-specific segmentation rendering and object annotation.
    \item Optional enhancement hooks for tracking and edge refinement.
    \item Lightweight geometric post steps compared to stator measurement-heavy path.
\end{itemize}

\subsection{File Post-Processing}
\begin{itemize}
    \item Per-bounding-box color analysis (blue/yellow/green scoring).
    \item Left-right layout rule validation with correctness message.
    \item Center-line drawing between selected file pair.
    \item UI toggle to enable/disable file color-layout validation at runtime.
\end{itemize}

\section{Deployment Decision and Current Status}

\begin{table}[H]
    \centering
    \begin{tabular}{lll}
        \toprule
        \textbf{Domain} & \textbf{Active Runtime Strategy} & \textbf{State} \\
        \midrule
        Stator & YOLO + UNet with measurement pipeline & Stable \\
        Chignon & YOLO + UNet via dedicated manager & Integrated \\
        File & UNet-ResNet18 + color/layout validation & Integrated, model lifecycle experimental \\
        \bottomrule
    \end{tabular}
    \caption{Current deployment state by domain.}
\end{table}

\section{Conclusion}

The report confirms the platform has matured into a domain-aware architecture with separate detection pipelines and post-processing methods. Stator and Chignon domains show strong model convergence, while File-domain YOLO training is currently non-viable. As a result, file deployment uses UNet-ResNet18 with dedicated color/layout validation logic until a successful YOLO retraining cycle is achieved.

\end{document}
